% !TEX TS-program = lualatex
% !TEX encoding = UTF-8

\documentclass[Session2026.tex]{subfiles}

\ifcsname preamble@file\endcsname
  \setcounter{page}{\getpagerefnumber{M-05_jeudi}}
\fi

\begin{document}
\addcontentsline{toc}{chapter}{Jeudi 9 juillet, Vierge Marie, Mère de Providence}


\newcommand{\frflex}[1]{#1} % TODO pouruqoi ça la ?

\bigtitle{Vierge Marie, Mère de Providence, à Vigiles}{Marie Mère de Providence}{Vigiles}

\rubric{On chante trois fois:}
\gscore{or_dlm}

\smallscore{or_dir}
\smalltitle{Psaume 3}
\psalm{3}{dir}

\newpage
\smalltitle{Invitatoire}
\gscore{inv_ave_maria}
\translation{\aa Je vous salue Marie, pleine de grâce, le Seigneur est avec vous.}
\psalm{94}{VLrepet}

\vspace{4mm}

\smalltitle{Hymne}
\gscore{hy_quem_terra}

\translation{Celui que terre, mer, astres
vénèrent, adorent, annoncent,\\
Celui qui régit ce triple monde,
Marie le porte caché dans son sein.\\
~\\
Celui que lune, soleil et toutes choses
servent en tout temps,\\
est porté par les entrailles d’une jeune vierge,
toute pénétrée de la grâce céleste.\\
~\\
La bienheureuse mère, par la grâce,
dans l’arche de son sein,\\
renferme l’Artisan suprême
qui tient le monde dans sa main.\\
~\\
Bienheureuse, à la parole d’un messager du ciel,
féconde par le Saint-Esprit,\\
et son sein donne au monde
le désiré des nations.}

\smalltitle{Premier nocturne}

\gscore{an_liberasti}
\translation{\aa Tu as racheté le sceptre de ton héritage.}
\smalltitle{Psaume 73}
\psalm{73}{8}
\smalltitle{Psaume 74}
\psalm{74}{8}

\gscore{an_tu_es_deus_qui_facis}
\translation{\aa Tu es le Dieu qui fais des merveilles.}
\smalltitle{Psaume 76}
\psalm{76}{6}

\gscore{an_inclinate_aurem}
\translation{\aa Prête l'oreille aux paroles de ma bouche.}
\smalltitle{Psaume 77}
\psalm{77-i}{1}
\smalltitle{Division}
\psalm{77-ii}{1}

\gscore{an_propitius_esto}
\translation{\aa Sois indulgent pour nos péchés, Seigneur.}
\psalm{78}{8}


\versiculus{Deus, in sancto via tua}{Quis Deus magnus sicut Deus noster}{Dieu, tes voies sont saintes}{Quel Dieu est grand comme notre Dieu?}

\newpage

\twocoltext{
	\vv Pater Noster...\\
	\vv Et ne nos indúcas in tentatiónem.\\
	\rr Sed líbera nos a malo.
}{
	\vv Notre Père...\\
	\vv Et ne nous laisse pas entrer en tentation.\\
	\rr Mais délivre-nous du mal.
}

\twocoltext{
	\vv Exáudi, Dómine Iesu Christe, preces servórum tuó\textit{rum},~\pscross{} et mise\textit{rére} \textbf{no}bis:~\psstar{} Qui cum Patre et Spíritu Sancto vivis et regnas in sǽcula sæculórum.\\
	\rr Amen.
}{
	\vv Seigneur Jésus-Christ, entends les prières de tes serviteurs, et aie pitité de nous, toi qui vis et règnes avec le Père et l'Esprit Saint dans les siècles des siècles.\\
	\rr Amen.
}

\vspace{4mm}

\smalltitle{Première lecture}

\twocoltext{
	\vv Iube, domne, benedícere.\\
	\vv Ille nos benedícat, qui sine fine vivit et regnat.\\
	\rr Amen.
}{
	\vv Maître, veuillez bénir.\\
	\vv Que celui qui vit et règne à jamais nous bénisse.\\
	\rr Amen.
}

\rubric{Après la lecture:}

\twocoltext{
	\vv Tu autem, Dómine, miserére nobis.\\
	\rr Deo grátias.
}{
	\vv Et toi, Seigneur, aie pitié de nous.\\
	\rr Nous rendons grâces à Dieu.
}

\gscore{re_sicut_cedrus}

\vspace{4mm}

\smalltitle{Deuxième lecture}

\twocoltext{
	\vv Iube, domne, benedícere.\\
	\vv Cuius festum cólimus, ipsa Virgo vírginum intercédat pro nobis ad Dóminum.\\
	\rr Amen.
}{
	\vv Maître, veuillez bénir.\\
	\vv Que celle dont nous célébrons la fête, la Vierge des vierges elle-même, intercède pour nous auprès du Seigneur.\\
	\rr Amen.
}

\newpage

\rubric{Après la lecture:}

\twocoltext{
	\vv Tu autem, Dómine, miserére nobis.\\
	\rr Deo grátias.
}{
	\vv Et toi, Seigneur, aie pitié de nous.\\
	\rr Nous rendons grâces à Dieu.
}

\gscore{re_gaude_maria}

\vspace{4mm}

\smalltitle{Troisième lecture}

\twocoltext{
	\vv Iube, domne, benedícere.\\
	\vv Ad societátem cívium supernórum perdúcat nos Rex Angelórum.\\
	\rr Amen.
}{
	\vv Maître, veuillez bénir.\\
	\vv Que le Roi des Anges nous fasse parvenir à la société des citoyens célestes.\\
	\rr Amen.
}

\rubric{Après la lecture:}

\twocoltext{
	\vv Tu autem, Dómine, miserére nobis.\\
	\rr Deo grátias.
}{
	\vv Et toi, Seigneur, aie pitié de nous.\\
	\rr Nous rendons grâces à Dieu.
}

\newpage

\gscore{re_ornatam_in_monilibus}

\newpage

\smalltitle{Deuxième nocturne}

\gscore{an_alleluia_feria5}
\smalltitle{Psaume 79}
\psalm{79}{6}
\smalltitle{Psaume 80}
\psalm{80}{6}
\smalltitle{Psaume 81}
\psalm{81}{6}
\smalltitle{Psaume 82}
\psalm{82}{6}
\smalltitle{Psaume 83}
\psalm{83}{6}
\smalltitle{Psaume 84}
\psalm{84}{6}

\gscore{an_alleluia_feria5} 

\smalltitle{Lecture brève}

\rubric{Après la lecture brève:}

\versiculus{O quam gloriósa est Mater}{Quæ cæli regem génuit}{Qu'elle est glorieuse, la Mère}{Qui a enfanté le roi des cieux}

\smalltitle{Intercession}
\smallscore{or_kyrie_vigilias_festivus}

\oratio{Omnípotens, sempitérne Deus, Ecclésiam tuam, déxtera poténtiæ tuæ, a cunctis prótege perículis,~\pscross{} et beáta María Matre Providéntiæ intercedénte,~\psstar{} fac eam præsénti gaudére prosperitáte et futúra. Per Dóminum.}{Dieu éternel et tout puissant, protège ton Église de tous les dangers par la force de ton bras, et par l'intercession de la bienheureuse Marie, Mère de Providence, fais-la se réjouir de sa prospérité présente et future. Par Jésus-Christ.}

\vspace{1mm}

\versiculus{Dóminus vobíscum}{Et cum spíritu tuo}{Le Seigneur soit avec vous}{Et avec votre esprit}

\vspace{1mm}
\gscore{or_benedicamus_vigilias_festis}

\versiculus{Fidélium ánimæ per misericórdiam Dei requiéscant in pace}{Amen}{Que par la miséricorde de Dieu, les âmes des fidèles trépassés reposent en paix}{Amen}

\bigtitle{Vierge Marie, Mère de Providence,\\à la Messe}{jeudi 9 juillet}{Messe}


\gscore{in_vultum_tuum}
\translation{\aa Tous les riches du peuple imploreront votre visage. On amènera au Roi des vierges à sa suite: ses proches vous seront présentées dans l’allégresse.\\
\vv D'heureuses paroles jaillissent de mon coeur quand je dis mes poèmes pour le roi.}

\rubric{Psalmodie de Tierce, comme au mercredi, p. \pageref{an_clamavi}.}

\smalltitle{Kyrie ad lib. VIII}
\gscore{ky_k8adlib}
\gscore{gr_adjutor_in_opportunitatibus}
\translation{\rr Secours au temps voulu, dans la détresse :
qu’ils mettent en vous leur espérance,
ceux qui vous connaissent :
car vous n’abandonnez pas ceux qui vous cherchent, Seigneur.\\
\vv Car le pauvre n'est pas oublié pour toujours: jamais ne périt l'espoir des malheureux. Lève-toi, Seigneur: qu'un mortel ne soit pas le plus fort.}

\gscore{al_domine_refugium}
\translation{\vv Le Seigneur fut pour nous un refuge, de génération en génération.}

\gscore{of_ave_maria}

\smalltitle{Sanctus X}
\gscore{ky_s10}
\smalltitle{Agnus X}
\gscore{ky_a10}
\gscore{co_laetabimur}
\translation{\aa Nous nous réjouirons de ton salut, nous serons glorifiés dans le nom du Seigneur notre Dieu.\\
\vv Que le Seigneur te réponde au jour de détresse, que le nom du Dieu de Jacob te défende.\\
\vv Du sanctuaire, qu'il t'envoie le secours, qu'il te soutienne des hauteurs de Sion.}


\bigtitle{Vierge Marie, Mère de Providence, à Sexte}{jeudi 9 juillet}{Sexte}

\smallscore{or_dia_ferialis}

\gscore{hy_rector_potens_memoriis}
\translation{\colored{M}aître puissant, Dieu Vérité, tu règles la marche du temps, tu formes l’aube en sa clarté, au midi tu donnes ses flammes.\\
\colored{E}teins le feu des dissensions, calme la fièvre du péché, apporte à nos corps la santé, à nos cœurs, la paix véritable.\\
\colored{E}xauce-nous, Père très bon, et toi, le Fils égal au Père, avec l’Esprit Consolateur, régnant pour les siècles des siècles.}

\rubric{Psalmodie de Sexte, comme au mardi, p. \pageref{an_qui_habitas}.}

\capitulum{Gal 5, 16-17}
{Spíritu ambuláte et concupiscéntiam carnis ne perfecéritis. / Caro enim
concupíscit advérsus Spíritum, Spíritus autem advérsus carnem ; † hæc
enim ínvicem adversántur, * ut non quæcúmque vultis illa faciátis.}
{Marchez sous la conduite de
l’Esprit Saint, et vous ne
risquerez pas de satisfaire les
convoitises de la chair. Car les
tendances de la chair s’opposent
à l’Esprit, et les tendances de
l’Esprit s’opposent à la chair. En
effet, il y a là un affrontement
qui vous empêche de faire tout
ce que vous voudriez.}

\versiculus{Bonus es tu, Dómine, et benefáciens}{Doce me iustificatiónes tuas}{Toi, tu es bon, Seigneur, tu fais du bien}{Apprends-moi tes commandements}

\rubric{Oraison des Vêpres, p. \pageref{0709_oratio}.}

\gscore{or_benedicamus_hm_ferialis}

\bigtitle{Vierge Marie, Mère de Providence, à None}{jeudi 9 juillet}{None}

\gscore{hy_rerum_deus_memoriis}
\translation{\colored{D}ieu fort, soutien de l’univers, qui es en toi sans changement, tu fixes dans leur succession les temps que le soleil mesure.\\
\colored{D}onne-nous un soir lumineux où la vie ne décline pas, où, pour fruit de la Sainte Mort, resplendira toujours la gloire.\\
\colored{E}xauce-nous, Père très bon, et toi, le Fils égal au Père, avec l’Esprit Consolateur, régnant pour les siècles des siècles.}

\rubric{Psalmodie de None, comme au mardi, p. \pageref{hy_rerum_deus_ferialis}.}

\capitulum{Gal 5, 22-23a.25}
{Fructus Spíritus est cáritas, gáudium, pax, longanímitas, * benígnitas,
bónitas, fides, mansuetúdo, continéntia. / Si vívimus Spíritu, Spíritu et
ambulémus.}
{Voici le fruit de l’Esprit :
amour, joie, paix, patience,
bonté, bienveillance, fidélité,
douceur et maîtrise de soi.
Puisque l’Esprit nous fait vivre,
marchons sous la conduite de
l’Esprit.}

\versiculus{Notam fac mihi, Dómine, viam in qua ámbulem}{Spíritus tuus bonus dedúcat me in terram rectam}{Montre-moi le chemin que je dois prendre}{Ton souffle est bienfaisant : qu’il me guide en un pays de plaines}

\oratio{
Da nobis orántibus, quǽsumus, Dómine,~\pscross{} ut patiéntiæ
Unigéniti tui sequámur exémpla,~\psstar{} et advérsa
patiéndi constántiam habeámus. Per Christum.
}{
Nous t’en prions, Seigneur,
que la passion de ton Fils
unique demeure devant nos
yeux et nous fortifie dans les
épreuves. Par le Christ.
}

\gscore{or_benedicamus_hm_ferialis}

\end{document}